\documentclass[12pt, letterpaper]{article}

% --- PREÁMBULO ---
\usepackage[utf8]{inputenc}
\usepackage[T1]{fontenc}
\usepackage[spanish]{babel}
\usepackage{geometry}
\geometry{top=2.5cm, bottom=2.5cm, left=2cm, right=2cm}
\usepackage{xcolor}
\usepackage{enumitem}

% --- FUENTE ---
\usepackage{helvet}
\renewcommand{\familydefault}{\sfdefault}

% =================================================================
% CONFIGURACIÓN PARA MOSTRAR / OCULTAR RESPUESTAS
% =================================================================
\newif\ifshowanswers

% ---> MODO PROFESOR (Muestra la franja y la respuesta justificada en azul).
% ---> MODO ESTUDIANTE (Pon un % al inicio de la línea para ocultarlas).
%\showanswerstrue
\showanswersfalse

% --- COMANDOS ESTILO "TAU" ---
% Se define con 3 argumentos: 
% #1: Opción correcta (letra)
% #2: Justificación del BIEN
% #3: Justificación del MAL (opcional)
\newcommand{\retroalimentacion}[3]{
	\ifshowanswers
	\vspace{0.2cm}
	\noindent\textcolor{gray}{\rule{\textwidth}{2pt}} % Franja gruesa simulando el libro
	\par\vspace{0.1cm}
	\noindent \textbf{\textcolor{blue}{BIEN: #1 \quad \textit{(#2)}}}
	
	% Verifica si el argumento 3 no está vacío
	\if\relax\detokenize{#3}\relax
	\else
	\par\vspace{0.1cm}
	\noindent \textbf{\textcolor{red}{MAL: \textit{#3}}}
	\fi
	\vspace{0.6cm}
	\else
	\vspace{1.2cm} % Espacio en blanco para el estudiante
	\fi
}

\begin{document}
	
	% --- ENCABEZADO ---
	\begin{center}
		{\Large \textbf{LICEO V.A.L. (Vida, Amor, Luz)}}\\[0.1cm]
		{\large DEPARTAMENTO DE CIENCIAS BÁSICAS}\\[0.1cm]
		\textbf{Taller de visto OK} \\[0.1cm]
		Ciencias 1 - TAU 7: Midamos la energía
	\end{center}
	
	\vspace{0.3cm}
	\noindent \textbf{Nombre:} \underline{\hspace{8cm}} \hspace{1cm}\textbf{Fecha:} \underline{\hspace{3cm}}
	
	\vspace{0.5cm}
	
	% --- PREGUNTA 1 ---
	\noindent \textbf{1.} La unidad de trabajo que se inventó el inglés James Joule hace unos 150 años se bautizó con su apellido. ¿A qué cantidad de esfuerzo físico equivale exactamente \textbf{un (1) julio}?
	\begin{enumerate}[label=\alph*)]
		\item A calentar un litro de agua en la estufa.
		\item A levantar un huevo de gallina (normal) desde el piso hasta la altura de la cintura.
		\item A subir un bulto de cemento hasta un quinto piso.
		\item A correr cien metros planos a toda velocidad.
	\end{enumerate}
	\retroalimentacion{b}{Ese es un trabajo fácil de ver y medir. Un huevo pesa más o menos la décima parte de un kilo (100 gramos), y al subirlo 1 metro hacemos 1 julio de trabajo.}{a, c y d. Esas opciones requieren cantidades inmensas de energía, muchísimo mayores a un solo julio.}
	
	% --- PREGUNTA 2 ---
	\noindent \textbf{2.} La regla matemática del TAU nos dice que "subir un kilo de peso a una altura de un metro equivale a efectuar 10 julios de trabajo". Sabiendo esto, si subes un morral de \textbf{10 kilos} por una escalera de \textbf{5 metros} de altura, efectúas un trabajo de:
	\begin{enumerate}[label=\alph*)]
		\item 15 julios.
		\item 50 julios.
		\item 500 julios.
		\item 5.000 julios.
	\end{enumerate}
	\retroalimentacion{c}{10 kilos equivalen a unos 100 huevos. Por cada metro que subes esos 10 kilos, haces un trabajo de 100 julios. Como son 5 metros: 100 x 5 = 500 julios.}{(10 kilos x 5 metros x 10 = 500).}
	
	% --- PREGUNTA 3 ---
	\noindent \textbf{3.} María va a sacar agua del fondo de un pozo que tiene una profundidad de 5 metros. Saca el balde lleno de agua (el agua pesa 7 kilos y el balde vacío pesa 1 kilo). ¿Cuánto trabajo total efectúa María?
	\begin{enumerate}[label=\alph*)]
		\item 50 julios.
		\item 400 julios.
		\item 350 julios.
		\item 500 julios.
	\end{enumerate}
	\retroalimentacion{b}{El peso total que levanta María es de 8 kilos (7 de agua + 1 del balde). 8 kilos x 5 metros x 10 = 400 julios de trabajo mecánico.}{c. El peso del balde también cuenta; María también tiene que subirlo.}
	
	% --- PREGUNTA 4 ---
	\noindent \textbf{4.} Cuando uno hace un esfuerzo físico fuerte, no solo efectúa trabajo mecánico (como mover el objeto). También realiza trabajos "internos", que son:
	\begin{enumerate}[label=\alph*)]
		\item Trabajo eléctrico y trabajo eólico.
		\item Trabajo bioquímico (respiración y circulación) y trabajo calórico (el cuerpo desprende calor).
		\item Trabajo de rozamiento y trabajo de pensamiento.
	\end{enumerate}
	\retroalimentacion{b}{Cuando haces ejercicio, la respiración y los latidos se aceleran (trabajo bioquímico) y el cuerpo desprende gran cantidad de calor.}{}
	
	% --- PREGUNTA 5 ---
	\noindent \textbf{5.} Científicamente, por cada 1 julio de trabajo externo (mecánico), el cuerpo gasta unos 3 julios en trabajo interno (bioquímico y calórico). Si recoges un objeto del piso y realizas un trabajo externo de \textbf{100 julios}, en total tu cuerpo ha efectuado un trabajo de:
	\begin{enumerate}[label=\alph*)]
		\item 100 julios.
		\item 300 julios.
		\item 400 julios.
	\end{enumerate}
	\retroalimentacion{c}{100 julios de trabajo externo (mecánico) más el triple de trabajo interno (300 julios). Total: 100 + 300 = 400 julios.}{a y b. Recuerda que siempre debes sumar ambos trabajos para saber cuánto gastó el organismo.}
	
	% --- PREGUNTA 6 ---
	\noindent \textbf{6.} ¿De qué se trataba la máquina o el "experimento mental" del señor Joule?
	\begin{enumerate}[label=\alph*)]
		\item Una piedra cayendo que movía un molinillo de paletas dentro de un recipiente cerrado con agua.
		\item Una estufa eléctrica calentando un litro de leche.
		\item Un globo que flotaba para medir el viento.
	\end{enumerate}
	\retroalimentacion{a}{La piedra al caer movía las paletas. El rozamiento agitaba el agua y la energía del movimiento se transformaba en calor.}{}
	
	% --- PREGUNTA 7 ---
	\noindent \textbf{7.} Con el experimento del molinillo, Joule logró descubrir el número exacto de trabajo necesario para calentar \textbf{un litro de agua en un grado centígrado}. Ese resultado fue:
	\begin{enumerate}[label=\alph*)]
		\item Exactamente 100 julios.
		\item 8.560 julios.
		\item Exactamente 4.280 julios.
		\item 1.000 julios.
	\end{enumerate}
	\retroalimentacion{c}{Ese es el valor exacto. Joule calculó que se necesitan exactamente 4.280 julios de trabajo para subir un grado a un litro de agua.}{}
	
	% --- PREGUNTA 8 ---
	\noindent \textbf{8.} Tienes un litro de agua a \textbf{20º} centígrados y lo pones al fogón para calentarlo hasta una temperatura de \textbf{30º} centígrados. ¿Cuánto trabajo calórico tuvo que hacer la estufa?
	\begin{enumerate}[label=\alph*)]
		\item 4.280 julios.
		\item 42.800 julios.
		\item 85.600 julios.
		\item 128.400 julios.
	\end{enumerate}
	\retroalimentacion{b}{La temperatura subió 10 grados (de 20º a 30º). Por cada grado son 4.280 julios. 10 $\times$ 4.280 = 42.800 julios.}{a. Ese valor sería si la calentaras un solo grado.}
	
	% --- PREGUNTA 9 ---
	\noindent \textbf{9.} Según el termómetro, a nivel del mar, el agua se congela y hierve exactamente a las siguientes temperaturas:
	\begin{enumerate}[label=\alph*)]
		\item Se congela a 10º centígrados y hierve a 90º centígrados.
		\item Se congela a 0º centígrados y hierve a 100º centígrados.
		\item Se congela a 15º centígrados y hierve a 160º centígrados.
	\end{enumerate}
	\retroalimentacion{b}{El agua se congela a los 0º y hierve al llegar a los 100º centígrados.}{}
	
	% --- PREGUNTA 10 ---
	\noindent \textbf{10.} Si estás en una ciudad alta en las cordilleras (como Bogotá), el agua ya no hierve a los 100º centígrados, sino a menor temperatura (ej. 92º). Esto se debe a que:
	\begin{enumerate}[label=\alph*)]
		\item En las ciudades altas los fogones calientan más lento por el frío.
		\item El agua de la montaña es más pura y hierve diferente.
		\item A mayor altura hay menor presión atmosférica, facilitando que el agua se evapore.
	\end{enumerate}
	\retroalimentacion{c}{Exacto. Esta técnica de medir la altura por la ebullición se llama "Hipsometría".}{}
	
\end{document}